\chapter{Number Systems}
\label{chap:number_systems}

\begin{tcolorbox}[enhanced, colback=boxnavybg, colframe=brandnavy, boxrule=1.2pt, arc=4pt, title={\Large\bfseries\sffamily Unit 4 Overview \& Learning Objectives}]
\sffamily\normalsize
This chapter establishes the mathematical principles of positional numeral systems used in computing. By the end of this unit, students will be able to:
\begin{itemize}[leftmargin=1.2em, itemsep=2pt, topsep=3pt]
    \item Understand the positional value formula and the role of the radix point, MSD, and LSD.
    \item Master binary, octal, decimal, and hexadecimal representations.
    \item Convert between any arbitrary base and the decimal number system.
    \item Execute rapid direct conversions between binary, octal, and hexadecimal using 3-bit and 4-bit grouping rules.
    \item Perform integer conversions via repeated division and fractional conversions via repeated multiplication.
\end{itemize}
\end{tcolorbox}

\section{Orientation: Positional Representation of Numbers}

A \textbf{number system} is a formal method of representing quantities using an agreed-upon set of symbols called \textbf{digits}. Modern computing is founded entirely on \textbf{positional number systems}, where the value contributed by any digit depends both on the symbol itself and its specific position within the number string.

\subsection{Base or Radix}
The \textbf{base} (or \textbf{radix}), denoted by $b$, defines the total number of distinct allowable digit symbols in that system. A base-$b$ system uses digits strictly ranging from $0$ up to $b-1$.
\begin{itemize}[leftmargin=1.2em]
    \item \textbf{Binary (Base 2):} Uses two digits: $0, 1$.
    \item \textbf{Octal (Base 8):} Uses eight digits: $0, 1, 2, 3, 4, 5, 6, 7$.
    \item \textbf{Decimal (Base 10):} Uses ten digits: $0, 1, 2, 3, 4, 5, 6, 7, 8, 9$.
    \item \textbf{Hexadecimal (Base 16):} Uses sixteen symbols: $0\text{--}9$ and $\text{A}\text{--}\text{F}$.
\end{itemize}

\subsection{The General Positional Evaluation Formula}
Any real number represented in a base-$b$ positional system can be evaluated into its equivalent decimal magnitude using the sum of products of each digit multiplied by its integer positional power of the base:
\begin{equation}
(d_n d_{n-1} \dots d_1 d_0 \,.\, d_{-1} d_{-2} \dots)_b = \sum_{i=-\infty}^{n} d_i \times b^i
\end{equation}
\begin{itemize}[leftmargin=1.2em]
    \item For integers, positional exponents increase positively from right to left starting at index $i = 0$ (the units position, with weight $b^0 = 1$).
    \item After the \textbf{radix point}, fractional exponents decrease negatively from left to right ($b^{-1}, b^{-2}, b^{-3}, \dots$).
    \item \textbf{MSD (Most Significant Digit):} The leftmost non-zero digit carrying the highest positional weight.
    \item \textbf{LSD (Least Significant Digit):} The rightmost digit carrying the lowest positional weight.
\end{itemize}

\section{The Four Key Number Systems in Computing}

\subsection{1. Binary (Base 2)}
Binary is the fundamental language of digital electronic computers because two-state electrical switches (transistors operating in saturation or cutoff) are physically simple, highly reliable, and resistant to electrical noise.
\begin{itemize}[leftmargin=1.2em]
    \item \textbf{Symbols:} $0, 1$.
    \item \textbf{Integer Weights:} $2^0 = 1, \; 2^1 = 2, \; 2^2 = 4, \; 2^3 = 8, \; 2^4 = 16, \; 2^5 = 32, \; 2^6 = 64, \; 2^7 = 128$.
    \item \textbf{Fractional Weights:} $2^{-1} = 0.5, \; 2^{-2} = 0.25, \; 2^{-3} = 0.125, \; 2^{-4} = 0.0625$.
\end{itemize}

\begin{example}[Evaluating a Binary Positional Number]
Convert $(1011.01)_2$ to its decimal equivalent:
\begin{align*}
(1011.01)_2 &= (1 \times 2^3) + (0 \times 2^2) + (1 \times 2^1) + (1 \times 2^0) + (0 \times 2^{-1}) + (1 \times 2^{-2}) \\
&= 8 + 0 + 2 + 1 + 0 + 0.25 \\
&= \mathbf{11.25_{10}}
\end{align*}
\end{example}

\begin{definition}[Binary Arithmetic Addition Rules]
\begin{align*}
0 + 0 &= 0 \\
0 + 1 &= 1 \\
1 + 0 &= 1 \\
1 + 1 &= 10_2 \quad (\text{Sum } 0 \text{ with a carry of } 1 \text{ to next column}) \\
1 + 1 + 1 &= 11_2 \quad (\text{Sum } 1 \text{ with a carry of } 1 \text{ to next column})
\end{align*}
\end{definition}

\subsection{2. Decimal (Base 10)}
Decimal is the universal human counting system. In computing, decimal is used at user interfaces for inputs and displayed reports. However, internally, computers do not store numbers in decimal. 

\begin{examalert}[Why Floating-Point Decimals Experience Rounding Errors]
In decimal, fractions like $1/10 = 0.1_{10}$ terminate cleanly. But in binary, because $10$ is not a pure power of $2$, the fraction $0.1_{10}$ becomes an \textbf{infinitely repeating binary fraction}:
\[
0.1_{10} = (0.0001100110011\dots)_2
\]
When stored in finite registers, the number must be truncated, resulting in tiny mathematical rounding discrepancies across repeated calculations.
\end{examalert}

\subsection{3. Octal (Base 8)}
Octal provides a convenient, human-friendly shorthand representation for binary strings. Because $8 = 2^3$, \textbf{exactly three binary bits correspond to one octal digit}.
\begin{itemize}[leftmargin=1.2em]
    \item \textbf{Digits Allowed:} $0, 1, 2, 3, 4, 5, 6, 7$ (the digits $8$ and $9$ are strictly illegal).
    \item \textbf{Weights:} $8^0 = 1, \; 8^1 = 8, \; 8^2 = 64, \; 8^3 = 512$.
    \item \textbf{Practical Application:} Traditional Unix/Linux file access permissions are expressed in octal (e.g., permission \texttt{755}, where octal digit $7_8 = 111_2$ represents read, write, and execute permissions).
\end{itemize}

\subsection{4. Hexadecimal (Base 16)}
Hexadecimal is the primary compact notation used across modern software engineering, computer architecture, and networking. Because $16 = 2^4$, \textbf{exactly four binary bits map to one hexadecimal digit}.
\begin{itemize}[leftmargin=1.2em]
    \item \textbf{Digits Allowed:} $0, 1, 2, 3, 4, 5, 6, 7, 8, 9$ and the letters $\text{A}=10, \text{B}=11, \text{C}=12, \text{D}=13, \text{E}=14, \text{F}=15$.
    \item \textbf{Weights:} $16^0 = 1, \; 16^1 = 16, \; 16^2 = 256, \; 16^3 = 4096$.
    \item \textbf{Syntax Conventions:} Often prefixed with \texttt{0x} in C/C++/Python (e.g., \texttt{0x2AF}) or suffixed with \texttt{h} in assembly (e.g., \texttt{2AFh}).
    \item \textbf{Applications:} Displaying memory addresses (e.g., \texttt{0x7FFF}), inspecting compiled machine code, defining MAC hardware addresses, and representing RGB colors in web design (e.g., \texttt{\#FF0000} represents pure red).
\end{itemize}

\section{Systematic Number Conversions}

\subsection{1. Converting from Any Base $b$ to Decimal}
Multiply each digit by its positional weight $b^i$ and compute the arithmetic sum:
\begin{example}[Hexadecimal to Decimal]
Convert $(2\text{AF})_{16}$ to decimal:
\begin{align*}
(2\text{AF})_{16} &= (2 \times 16^2) + (10 \times 16^1) + (15 \times 16^0) \\
&= (2 \times 256) + (10 \times 16) + (15 \times 1) \\
&= 512 + 160 + 15 = \mathbf{687_{10}}
\end{align*}
\end{example}

\subsection{2. Converting Decimal Integers to Any Base $b$ (Repeated Division)}
Divide the decimal integer repeatedly by the target base $b$, recording the integer quotient and the remainder at each step. Terminate when the quotient becomes $0$. The final answer is obtained by reading the remainders \textbf{from bottom (last remainder = MSD) to top (first remainder = LSD)}.

\begin{example}[Decimal $156_{10}$ to Binary]
\begin{align*}
156 \div 2 &= 78 \quad \text{remainder } 0 \quad (\text{LSD}) \\
78 \div 2  &= 39 \quad \text{remainder } 0 \\
39 \div 2  &= 19 \quad \text{remainder } 1 \\
19 \div 2  &= 9 \quad \text{remainder } 1 \\
9 \div 2   &= 4 \quad \text{remainder } 1 \\
4 \div 2   &= 2 \quad \text{remainder } 0 \\
2 \div 2   &= 1 \quad \text{remainder } 0 \\
1 \div 2   &= 0 \quad \text{remainder } 1 \quad (\text{MSD}) \quad \implies \mathbf{(10011100)_2}
\end{align*}
\end{example}

\subsection{3. Converting Decimal Fractions to Any Base $b$ (Repeated Multiplication)}
Multiply the fractional portion repeatedly by the target base $b$. Record the resulting integer part at each step and continue multiplying only the remaining fractional component. Read the integer parts \textbf{from top (first generated = MSD) to bottom (last generated = LSD)}.

\subsection{4. Direct Grouping Conversions (Fast Method)}
Because $8 = 2^3$ and $16 = 2^4$, conversions between binary, octal, and hexadecimal bypass decimal arithmetic entirely:
\begin{itemize}[leftmargin=1.2em]
    \item \textbf{Binary $\rightarrow$ Octal:} Partition bits into groups of \textbf{three} starting from the radix point (adding leading or trailing zeros if needed):
    \[
    \underbrace{010}_{2}\;\underbrace{011}_{3}\;\underbrace{100}_{4}\;_2 \;\implies\; \mathbf{(234)_8}
    \]
    \item \textbf{Binary $\rightarrow$ Hexadecimal:} Partition bits into groups of \textbf{four} starting from the radix point:
    \[
    \underbrace{1001}_{9}\;\underbrace{1100}_{12=\text{C}}\;_2 \;\implies\; \mathbf{(9\text{C})_{16}}
    \]
    \item \textbf{Octal $\leftrightarrow$ Hexadecimal:} Convert each octal digit to 3 binary bits, then regroup the binary sequence into 4-bit blocks to yield the hexadecimal equivalent.
\end{itemize}

\begin{example}[Integrated Equivalency Verification]
Notice that:
\[
\mathbf{156_{10}} \;=\; \mathbf{(10011100)_2} \;=\; \mathbf{(234)_8} \;=\; \mathbf{(9\text{C})_{16}}
\]
All four representations denote the exact same mathematical quantity.
\end{example}

\section{Chapter Review \& Key Takeaways}
\begin{tcolorbox}[enhanced, colback=boxgreenbg, colframe=boxgreenborder, boxrule=1pt, arc=4pt, title={\bfseries\sffamily Summary Checklist}]
\sffamily\small
\begin{itemize}[leftmargin=1.2em, itemsep=2pt]
    \item A positional number system evaluates value through $\sum (d_i \times b^i)$.
    \item Binary is base 2 ($0,1$); Octal is base 8 ($0\text{--}7$); Decimal is base 10 ($0\text{--}9$); Hexadecimal is base 16 ($0\text{--}9, \text{A}\text{--}\text{F}$).
    \item Octal digits map directly to 3 binary bits; Hexadecimal digits map directly to 4 binary bits.
    \item Decimal integers convert via repeated division (read remainders bottom-to-top); fractions convert via repeated multiplication (read integers top-to-bottom).
    \item An unsigned 8-bit byte can represent decimal values from $0$ up to $255$ ($11111111_2 = \text{FF}_{16}$).
\end{itemize}
\end{tcolorbox}
